\scanpage{72}
\mspage{71}
\begin{center}- 35 -\end{center}
\begin{flushright}n°357\end{flushright}
nous énonçons sous forme de Lemme :

LEMME 3.- \uline{Soient} $\lambda$, $\lambda'$ \uline{deux mesures composables sur} $G$ ; \uline{soit} $\psi$ \uline{un caractère (représentation continue) de} $G$ \uline{dans} $\mathbb C^*$ ; \uline{alors} $\psi,\lambda$ \uline{et} $\psi.\lambda'$ \uline{sont composables, et on a} :
\[
(\psi.\lambda)*(\psi.\lambda')=\psi.(\lambda*\lambda').
\]
On peut exprimer les résultats ci-dessus en disant que, par la loi $(\lambda,f)\to(\chi^{1/q}.\lambda)*f$, l'algèbre $\mathcal M_b^1(G)$ opère continument sur $L^p(E;\gamma)$.

PROPOSITION 9.- \uline{Soient} $\lambda$ \uline{une mesure bornée sur} $G$, $f\in\mathcal L^p$, $g\in\mathcal L^q$ ; \uline{alors on a} :
\[
\tag{14}
\int g.(\chi^{1/q}\lambda*f)\,d\gamma
=\int(\chi^{1/p}\lambda^\vee*g).f\,d\gamma,
\]
\uline{où} $\lambda^\vee$ \uline{est la mesure donnée par} $d\lambda^\vee(x)=d\lambda(x^{-1})$, \uline{c'est-à-dire l'image de} $\lambda$ \uline{par} $x\to x^{-1}$.

Observons d'abord qu'on a, pour $F\in\mathcal K(E)$ :
\[
\begin{aligned}
\int F(xz)\chi(x)^{1/q}\,d\lambda(x)
&=\int F(x^{-1}z)\chi(x)^{-1/q}\,d\lambda^\vee(x)=\\
&=\int F(x^{-1}z)\chi(x)^{1/p}.\chi(x^{-1})\,d\lambda^\vee(x)
=\int((\unsure{\chi|\lambda}^{1/p}\lambda^\vee)*F)(z).
\end{aligned}
\]
Dans ces conditions, si on applique au premier membre de la formule (2), dans la prop.2, le théorème de Lebesgue-Fubini, après y avoir remplacé $\mu$ par $\chi^{1/q}\lambda$, on obtient (14) avec $F$ au lieu de $g$ ; cela démontre (14) pour $g\in\mathcal K(E)$. Comme les deux membres dépendent continument de $g\in\mathcal L^p$ en vertu de la prop.8, et que $\mathcal K(E)$ est partout dense dans $\mathcal L^p$, la conclusion s'ensuit.

[Le rédacteur (resp. le lecteur) s'aperçoit avec plaisir (resp. chagrin) que la belle Prop.6 de la rédaction antérieure, (p.57), avec ses succulentes conditions $\mathrm{DM}_{\mathrm I}$, $\mathrm{DM}_{\mathrm {II}}$, $\mathrm{DM}_{\mathrm {III}}$, $\mathrm{DM}_{\mathrm {IV}}$, dont Bourbaki semblait si fier,a disparu sans laisser de trace, vu qu'il n'en a eu nul besoin pour la convolution dans les $L^p$, non plus que du n° hérétique du §1, ni de l'intégration vectorielle, ni des hypocons, etc. Tremblant de crainte à l'idée de déplaire à son illus-
